% HAND-WRITTEN fragment -- NOT generated by maestro2tex. Input from
% AnalogChapter.tex immediately after the assembled-channel fragments.
%
% Unlike the other note_meas_* fragments, part of this fabric EXISTS in the
% schematic: the converter's 16:1 input multiplexer and its 4-bit select pin
% (slot 0 = transimpedance output, 1 = common mode, 2 = commanded potential,
% 3 = reserved for a buffered RE tap and floating, 4 = analog test point,
% 5 = ground). What does not exist: the RE buffer (F-5c), the common-mode
% buffer (the S7.5 finding), the test-port mux inputs (F-3), the electrode
% isolation switches and the on-chip gain reference resistor.
% Numbers: bit-9 inversion and 88.5 nA/LSB at N = 0 (sss:pixtop-adc); 0.3 nV
% zero spread over 320 codes and 0.825 nA/LSB at N = 319 (sss:pixtop-zero);
% 0.731 V / 16 mV collapse of the common-mode slot (sss:pixtop-open).
% Library cell names are kept out of the rendered prose.

\subsubsection{Measuring this block on chip}

The channel's own converter multiplexer carries this sequence; only the last
step needs the tester.
\begin{enumerate}[nosep]
\item Invert bit~9 of the raw converter bus before anything else reads it. Every
  step below is that number.
\item Select the transimpedance output (slot 0) with the working electrode
  isolated and take the current zero at a \emph{high-gain} code. The zero is
  gain-independent to \SI{0.3}{\nano\volt} over all 320 taps, so one
  measurement calibrates every code --- but at the bottom code one LSB is
  \SI{88.5}{\nano\ampere} against \SI{34}{\nano\ampere} of offset and the
  measurement is quantisation noise.
\item Do not take that zero through the common-mode slot until that tap is
  buffered: sampling it pulls the common mode to \SI{0.731}{\volt} and the
  output to \SI{16}{\milli\volt} at every conversion, and the channel never
  settles.
\item Read the commanded potential at slot 2 to separate a DAC fault from an
  amplifier fault. A1's own contribution is not observable: slot 3 is reserved
  for a buffered reference-electrode tap that does not exist, and the raw
  electrode must not be sampled.
\item Absolute potential scale and transimpedance gain still need the
  production tester at the reference-DAC pad and the working-electrode pad
  (Sections~\ref{ss:dacr2r12} and~\ref{ss:tiarprog}). The in-loop feedback
  resistance of Section~\ref{sss:pixtop-cv} has no on-die counterpart --- it is
  read from a simulation probe.
\end{enumerate}
